\newpage
\subsection{\textcolor{Salmon}{Ergebnisse Nutzerbefragung}}
\textcolor{blue}{\textit{einleitende Worte}} \\

\noindent
\underline{Fehlerbetrachtung}\\
- Fehlerhafte Daten\\

% Boxplots
\noindent
\underline{s. \autoref{fig:pruemper_all}}\\
- zeigt alle Antworten der Fragen in einem Diagramm\\
- Tendenz zu neu besser Erkennbar (erst in Diskussion?)\\
\textcolor{red}{- Kurzbeschreibung: Erklärung, wie Boxplots aufgebaut sind und funktionieren\\
-> 2 Sätze} % hatten wir potenziell in Mathe 2 (beschreibende Statistik)

\noindent
\underline{s. \autoref{tab:Nutzerbefragung:mittelwerte}}\\
- Mittelwerte, Median, Standardabweichung\\
- Differenzen

\begin{table}[H]
\begin{center}
\begin{tabular}{|l|c|c|c|c|c|}
\hline
Frage                                   & 1      & 2      & 3        & 4      & 5      \\ \hline
$\varnothing_{neu} - \varnothing_{alt}$ & 1,73   & 1,18   & 0,91     & 1,73   & 1,55   \\ \hline
T-Test                                  & 0,0195 & 0,0949 & 0,138729 & 0,0487 & 0,0231 \\ \hline
\end{tabular}
\begin{tabular}{|l|c|c|c|c|c|}
\hline
Frage                                   & 1      & 2      & 3        & 4      & 5      \\ \hline
$\varnothing_{neu} - \varnothing_{alt}$ & 1,73   & 1,18   & 0,91     & 1,73   & 1,55   \\ \hline
T-Test                                  & 0,0195 & 0,0949 & 0,138729 & 0,0487 & 0,0231 \\ \hline
\end{tabular}
\caption{Nutzerbefragung - Differenz der Antwortmittelwerte und T-Test-Ergebnisse (Ergebnisse von Weiterverarbeitung der Daten aus \autoref{tab:Nutzerbefragung:rohe_Daten})}
\label{tab:Nutzerbefragung:Differenz+ttest}
\end{center}
\end{table}

\noindent
\underline{s. \autoref{tab:Nutzerbefragung:Differenz+ttest}}\\
- Excel-Formel gleiche wie in \autoref{Bsp:ttest}\\
- volle Excel-Daten in \autoref{fig:Usability_Test:rohe_Daten}
